\section{Rollup to the Z-Chain}
\label{sec:rollup}

Per-validator re-execution of every trade is a ceiling: a chain that re-runs each match on
each validator cannot finalize trades faster than one validator can re-run them. We break
this ceiling by rolling D-Chain trade batches up to the Z-Chain (the Lux ZKVM), which
verifies \emph{one proof per $K$ trades} instead of re-executing each trade. This section
gives the two-stage rollup, the measured verification rates, and an honest separation of
settlement \emph{finality} from matching \emph{throughput}.

\subsection{Batches and proofs}

The D-Chain produces, per market, a batch of $K$ fills together with the batch's execution
root. A proof attests that the batch is a valid extension of the prior root. The Z-Chain
verifies the proof and finalizes all $K$ trades in the batch at once. The effective
settlement-finality rate is
\begin{equation}
  \text{settled-tx/s (finality)} \;=\; K \times V \times M,
  \label{eq:rollup}
\end{equation}
where $V$ is the per-core proof-verification rate and $M$ the number of markets verified in
parallel. Equation~\eqref{eq:rollup} is a finality-axis relation: it counts trades finalized
per second by verification, which is independent of how fast trades are matched.

\subsection{Two stages: attestation then validity}

\begin{description}
  \item[P3Q (venue-trust attestation).] The D-Chain signs the batch with an ML-DSA-65
  (FIPS-204) signature; the Z-Chain verifies the signature at the P3Q precompile
  (\texttt{0x012205}). This is cheap to produce and to verify, and it is sound under the
  assumption that the single-operator venue is honest---appropriate for the current
  single-venue deployment.
  \item[\texttt{starkfri} (trustless validity).] The D-Chain produces a STARK/FRI validity
  proof that the batch transition is correct; the Z-Chain verifies it with no trust in the
  venue. This is the path to a decentralized D-Chain. The verifier is a real post-quantum
  STARK/FRI implementation (cSHAKE256 Fiat-Shamir transcript over the Goldilocks field), not
  a placeholder, and carries formal-verification artifacts for its on-chain envelope.
\end{description}

The two are complementary: P3Q gives cheap finality now under venue trust;
\texttt{starkfri} gives trustless finality as the venue decentralizes. The carried-fills
block layout (Section~\ref{sec:determinism}) reserves a signature field precisely so this
transition needs no further wire change.

\subsection{Measured verification rates}

\begin{table}[H]
\centering
\renewcommand{\arraystretch}{1.25}
\begin{tabular}{l l r l}
\toprule
\textbf{Stage} & \textbf{Primitive} & \textbf{Verifies/s ($V$)} & \textbf{Trust model} \\
\midrule
P3Q       & ML-DSA-65 attestation & 4{,}714 & venue-trust \\
\texttt{starkfri} & STARK/FRI validity proof & 986 & trustless \\
\bottomrule
\end{tabular}
\caption{Measured single-core verification rates (Apple M1 Max), through the full precompile
dispatch. P3Q is the real FIPS-204 verify; \texttt{starkfri} is the real STARK/FRI verifier.
\textbf{The \texttt{starkfri} figure is for a toy AIR} (degree 16, blowup 8, 16 queries); a
production DEX-batch AIR has many more constraints and verifies slower. We report 986/s as a
ceiling for that proof shape, not as a production-batch number; the production-batch AIR is
scaffolded and not yet measured.}
\label{tab:verify}
\end{table}

\subsection{Settlement finality at fleet scale}

Substituting the measured $V$ into Equation~\eqref{eq:rollup} shows the rollup clears
$10^{9}$ settled-tx/s on the \emph{finality} axis with realistic batch sizes and modest
parallelism:

\begin{table}[H]
\centering
\renewcommand{\arraystretch}{1.25}
\begin{tabular}{l r r r r}
\toprule
\textbf{Primitive} & \textbf{$K$} & \textbf{$M$} & \textbf{cores} & \textbf{settled-tx/s} \\
\midrule
P3Q       & 10{,}000  & 22 & 1  & $1.04\times10^{9}$ \\
P3Q       & 100{,}000 & 3  & 1  & $1.41\times10^{9}$ \\
\texttt{starkfri} & 10{,}000  & 22 & 10 & $2.17\times10^{9}$ \\
\bottomrule
\end{tabular}
\caption{Settlement-finality throughput from Equation~\eqref{eq:rollup} using the measured
$V$ of Table~\ref{tab:verify}. These combine a measured per-proof verify rate with a stated
batch size and parallelism; the verify rate is measured, the batch size and core count are
configuration. The \texttt{starkfri} row uses the toy-AIR rate and is therefore an
upper-bound shape, not a production-batch figure.}
\label{tab:finality}
\end{table}

\subsection{Finality is not matching}

We are explicit about what these numbers do and do not say. The rollup finalizes $K$ trades
per verified proof, so its rate (Equation~\eqref{eq:rollup}) is a \emph{settlement-finality}
rate---it counts already-matched trades being finalized. It is \emph{not} the rate at which
new trades are matched, which is the GPU-fleet-bound figure of Section~\ref{sec:gpu-matching}
($\sim$100--250\,M matches/s per node, aggregated across a fleet for $10^{9}$ matches/s).
Conflating the two is exactly the error this paper retracts. The three axes---matching
(per-book parallel), settlement (market-sharded Block-STM), and finality (rollup)---are
independent and independently measured.

\subsection{Private order flow}

Encrypted order flow can be carried inside a batch using fully homomorphic encryption, so
that order contents remain private while still being matched and proven. We note the honest
limit: homomorphic bootstrapping runs at single to low-tens of operations per second per
accelerator, orders of magnitude below the attestation rate, so the private path is bounded
by FHE bootstrap cost rather than by verification. Private order flow is therefore a distinct
performance regime from the public path and is governed by FHE throughput.
